\section*{B.1 附加误差统计}
\addcontentsline{toc}{section}{B.1 附加误差统计}
\begin{table}[H]
\centering
\caption{不同噪声水平下 RMSE（m）}
\begin{tabular}{cccc}
\toprule
噪声等级 & RSSI-Only & KF-Baseline & Proposed \\
\midrule
低噪声 & 3.21 & 2.45 & 1.78 \\
中噪声 & 5.08 & 3.62 & 2.54 \\
高噪声 & 7.44 & 5.27 & 3.69 \\
\bottomrule
\end{tabular}
\end{table}

\section*{B.2 信标失效率实验}
\addcontentsline{toc}{section}{B.2 信标失效率实验}
当失效率为 10\%、20\%、30\% 时，本文方法 RMSE 分别为 2.31 m、2.72 m、3.05 m，表现出较好的退化平稳性。

\section*{B.3 计算开销实验}
\addcontentsline{toc}{section}{B.3 计算开销实验}
在窗口长度 $w=8$、最大迭代 10 次配置下，单帧平均耗时约 2.4 ms，满足 1 Hz 实时定位要求。

\section*{B.4 复现说明}
\addcontentsline{toc}{section}{B.4 复现说明}
实验复现需要固定随机种子，并记录信标布局、噪声参数、窗口长度和正则系数。建议保留每次运行的日志文件以便结果追踪。
